\section{Factorized observable, kinematics, and hard-function definition}
\label{sec:factorized-observable}

\subsection{Hadronic observable and perturbative expansion}

We consider
\begin{equation}
 A(P_A,S_A)+B(P_B,S_B)\longrightarrow h(P_h,S_h)+X,
 \label{eq:hadronic-process}
\end{equation}
with hadron masses neglected in the short-distance coefficient.  The
single-particle invariant cross section is
\begin{equation}
 \mathcal I_h^{\mathsf P}(P_h)
 :=E_h\frac{\dd^3\sigma^{\mathsf P}}{\dd^3\bm P_h},
 \qquad
 \mathsf P\in\{\mathsf{UU},\mathsf{LL},\mathsf{TT}\}.
 \label{eq:invariant-observable}
\end{equation}
Hadron $B$ is unpolarized; the two letters label the polarization of
hadron $A$ and of the observed hadron $h$, respectively.

At leading twist, the factorized cross section may be written as
\begin{equation}
 \mathcal I_h^{\mathsf P}
 =\sum_{a,b,c}\int\dd x_a\,\dd x_b\,\dd z_h\;
 \mathcal J_h\,
 \mathscr F_{ab\to c}^{\mathsf P}(x_a,x_b,z_h;\mu_F)
 \frac{1}{x_ax_bz_h^2}\,
 \mathcal H_{ab\to c}^{\mathsf P}
 (s,t,u;\mu_R,\mu_F,\eps).
 \label{eq:factorized-cross-section}
\end{equation}
Here $\mathcal J_h$ contains the remaining Jacobian associated with the
chosen hadronic variables after the displayed fraction monomial has been
removed.  This notation is equivalent to the conventional
single-inclusive formula of Ref.~\cite{Aversa:1988vb}; it makes explicit the
factor whose cancellation is tested when the hard coefficient is extracted.
The distribution products $\mathscr F^{\mathsf P}$ are given below.

We keep the QCD coupling normalization explicit,
\begin{equation}
 \mathcal H^{\mathsf P}
 =\alpha_s^2(\mu_R)H^{\mathsf P,(0)}
 +\alpha_s^3(\mu_R)H^{\mathsf P,(1)}
 +\alpha_s^4(\mu_R)H^{\mathsf P,(2)}+\cdots .
 \label{eq:hard-perturbative-expansion}
\end{equation}
Factors of $2\pi$, the flux, initial-state averages, and the
$\overline{\mathrm{MS}}$ scale factor are part of the stated normalization
of each $H^{(n)}$; they are not absorbed into a second definition of the
coupling.  The present paper develops the exact real-emission and
master-integral ingredients.  A complete fixed-order hard function additionally
requires the virtual, ultraviolet-renormalization, and collinear-factorization
terms in this same convention.

\subsection{Hadronic and partonic kinematics}

Define
\begin{equation}
 S=(P_A+P_B)^2,
 \qquad T=(P_A-P_h)^2,
 \qquad U=(P_B-P_h)^2,
 \label{eq:hadronic-invariants}
\end{equation}
and the collinear map
\begin{equation}
 k_a=x_aP_A,
 \qquad k_b=x_bP_B,
 \qquad P_h=z_hk_c,
 \qquad 0<x_a,x_b,z_h<1.
 \label{eq:collinear-map}
\end{equation}
The partonic invariants are
\begin{equation}
 s=(k_a+k_b)^2,
 \qquad t=(k_a-k_c)^2,
 \qquad u=(k_b-k_c)^2,
 \label{eq:partonic-invariants}
\end{equation}
so that
\begin{equation}
 s=x_ax_bS,
 \qquad t=\frac{x_a}{z_h}T,
 \qquad u=\frac{x_b}{z_h}U.
 \label{eq:hadronic-partonic-map}
\end{equation}
The total recoil momentum and its invariant mass are
\begin{equation}
 q=k_a+k_b-k_c,
 \qquad Q^2=q^2=s+t+u.
 \label{eq:recoil-invariant}
\end{equation}
For real emission we work in the chamber
\begin{equation}
 s>0,\quad t<0,\quad u<0,\quad Q^2>0,\quad
 s+t>0,\quad s+u>0,\quad t+u<0.
 \label{eq:physical-chamber}
\end{equation}
All noninteger powers of positive quantities use the real logarithm in this
chamber.

The single-inclusive variables are
\begin{equation}
 v=1+\frac{t}{s},
 \qquad w=-\frac{u}{s+t},
 \qquad x=1-w,
 \label{eq:vw-definitions}
\end{equation}
with inverse map
\begin{equation}
 t=-s(1-v),
 \qquad u=-svw=-sv(1-x),
 \qquad Q^2=sv(1-w)=svx.
 \label{eq:vw-inverse}
\end{equation}
Thus $s>0$, $0<v,w<1$, and the real-emission endpoint is
\begin{equation}
 w\to1^-,\qquad x\to0^+,
 \qquad v\ \text{fixed}.
 \label{eq:real-endpoint}
\end{equation}

\subsection{Leading-twist polarization projections}

For each collinear direction $r\in\{A,B,h\}$, let
$n_r^2=\bar n_r^2=0$ and $n_r\cdot\bar n_r=1$.  The physical transverse
spins obey
\begin{equation}
 S_{AT}\cdot n_A=S_{AT}\cdot\bar n_A=0,
 \qquad
 S_{hT}\cdot n_h=S_{hT}\cdot\bar n_h=0.
 \label{eq:transverse-spin-conditions}
\end{equation}
They are hadronic vectors; the partonic polarization information enters
through the leading-twist correlators.

For $\chi=q,\bar q$, define $\eta_q=1$ and
$\eta_{\bar q}=-1$.  With
$i\sigma^{ab}=-[\not\!a,\not\!b]/2$, the distribution and fragmentation
correlators are
\begin{align}
 \Phi_{\chi/H}(x;\lambda_H,S_{HT})
 =\frac12\bigl[&f_1^{\chi/H}(x)\not\!n_H
 +\eta_\chi\lambda_Hg_{1L}^{\chi/H}(x)
   \gamma_5\not\!n_H
 +h_1^{\chi/H}(x)i\sigma^{n_HS_{HT}}\gamma_5\bigr],
 \label{eq:pdf-correlator}\\
 \Delta_{h/\chi}(z_h;\lambda_h,S_{hT})
 =\frac12\bigl[&D_1^{h/\chi}(z_h)\not\!n_h
 +\eta_\chi\lambda_hG_{1L}^{h/\chi}(z_h)
   \gamma_5\not\!n_h
 +H_1^{h/\chi}(z_h)i\sigma^{n_hS_{hT}}\gamma_5\bigr].
 \label{eq:ff-correlator}
\end{align}
The antiquark convention therefore changes the helicity term but not the
unpolarized or transversity term.  These are the standard collinear
leading-twist structures~\cite{Bacchetta:2006tn}.  The BMHV definition of
$\gamma_5$ is given in Appendix~\ref{app:conventions}.

The three distribution products in Eq.~\eqref{eq:factorized-cross-section}
are
\begin{align}
 \mathscr F^{\mathsf{UU}}_{ab\to c}
 &=f_1^{a/A}(x_a)f_1^{b/B}(x_b)D_1^{h/c}(z_h),
 \notag\\
 \mathscr F^{\mathsf{LL}}_{ab\to c}
 &=g_{1L}^{a/A}(x_a)f_1^{b/B}(x_b)G_{1L}^{h/c}(z_h),
 \notag\\
 \mathscr F^{\mathsf{TT}}_{ab\to c}
 &=h_1^{a/A}(x_a)f_1^{b/B}(x_b)H_1^{h/c}(z_h).
 \label{eq:polarized-distribution-products}
\end{align}
The LL hard part multiplies $\lambda_A\lambda_h$.  The TT hard tensor is
contracted with $S_{AT}^\mu S_{hT}^\nu$ and then decomposed into independent
physical azimuthal structures.  Each scalar coefficient in that decomposition
is treated as a separate hard function.

\subsection{Analytic content of the hard function}

After the distribution product, the convolution Jacobian, and the displayed
fraction monomial in Eq.~\eqref{eq:factorized-cross-section} are removed, the
hard function depends only on the partonic invariants, scales, color factors,
and $\eps$.  At fixed $(s,t,u)$, it therefore satisfies
\begin{equation}
 \frac{\partial H^{\mathsf P,(n)}}{\partial x_a}
 =\frac{\partial H^{\mathsf P,(n)}}{\partial x_b}
 =\frac{\partial H^{\mathsf P,(n)}}{\partial z_h}=0.
 \label{eq:fraction-independence}
\end{equation}
For UU and LL it is also independent of every auxiliary light-cone and
hadronic basis vector.  For TT the same statement applies to each scalar
coefficient after projection onto the physical azimuthal basis.  A common
scalar factor may be removed temporarily during exact algebra, but it is
restored before endpoint analysis; in particular, a factor
$(Q^2)^{-\eps}=(sv)^{-\eps}(1-w)^{-\eps}$ cannot be discarded.

Before expanding in $\eps$, write the real-emission endpoint terms as
\begin{equation}
 H_{\mathrm{real}}^{\mathsf P,(n)}(v,w,\eps)
 =H_{\mathrm{reg}}^{\mathsf P,(n)}(v,w,\eps)
 +\sum_{r=1}^{N_{\mathrm{end}}}
 C_r^{\mathsf P,(n)}(v,\eps)(1-w)^{-1-a_r\eps},
 \qquad a_r>0.
 \label{eq:unexpanded-endpoint-form}
\end{equation}
The branches $a_r$ and coefficients $C_r$ are extracted from the exact
reconstructed expression.  Defining
\begin{equation}
 \mathcal D_m(w)=\left[\frac{\log^m(1-w)}{1-w}\right]_+,
 \label{eq:plus-distributions}
\end{equation}
meromorphic continuation gives
\begin{equation}
 (1-w)^{-1-a\eps}
 =-\frac{\delta(1-w)}{a\eps}
 +\sum_{m=0}^{\infty}\frac{(-a\eps)^m}{m!}\mathcal D_m(w).
 \label{eq:endpoint-distribution-identity}
\end{equation}
Only after Eq.~\eqref{eq:unexpanded-endpoint-form} is established do we form
the fixed-order decomposition
\begin{equation}
 H^{\mathsf P,(n)}
 =H_\delta^{\mathsf P,(n)}(v)\delta(1-w)
 +\sum_{m=0}^{m_{\max}}H_m^{\mathsf P,(n)}(v)\mathcal D_m(w)
 +H_{\mathrm{reg}}^{\mathsf P,(n)}(v,w).
 \label{eq:hard-distributional-form}
\end{equation}
The coefficients in Eq.~\eqref{eq:hard-distributional-form} are exact
analytic functions.  Numerical values at chosen points can test them but do
not define them.
